\subsubsection{Продвинутые задания}

\begin{taskbox}[Плавная смена цвета этапов миссии]
Требуется сопоставить цвета состояниям миссии (подготовка / полёт / посадка / ожидание):
\begin{itemize}
    \item обновлять все индексы \texttt{0..28} без блокировок \texttt{sleep()};
    \item реализовать смену через функции стадий и отложенные вызовы.
\end{itemize}

\textbf{Ожидаемый результат:} плавная индикация переходов между этапами без пропусков и блокировок.
\end{taskbox}

\begin{taskbox}[Динамический радиус окружности по RC7]
Требуется считывать \texttt{Sensors.rc()} и управлять радиусом:
\begin{itemize}
    \item нормировать значение канала к \texttt{[0, 1]};
    \item вычислять радиус \(r = 0.2 + 0.8 \cdot \text{norm}\);
    \item периодически обновлять точку окружности.
\end{itemize}

\textbf{Ожидаемый результат:} радиус плавно меняется от \texttt{0.2} до \texttt{1.0} по RC7, траектория стабильна.
\end{taskbox}

\begin{taskbox}[Ограничение скорости изменения угла \texttt{angle}]
Требуется реализовать \texttt{advanceAngle(step)}:
\begin{itemize}
    \item с ограничением \texttt{step} диапазоном \texttt{[-maxStep, maxStep]};
    \item применять его при каждом обновлении параметра окружности.
\end{itemize}

\textbf{Ожидаемый результат:} угол изменяется с ограниченной скоростью, резкие скачки отсутствуют.
\end{taskbox}

\begin{taskbox}[Перезапуск миссии после \texttt{Ev.COPTER\_LANDED}]
Требуется по событию выполнить \texttt{Timer.callLater(3, fn)} и инициировать новый взлёт, исключая блокирующие задержки.

\textbf{Ожидаемый результат:} автоматический повтор миссии после посадки с паузой, без блокировок.
\end{taskbox}

\begin{taskbox}[«Морзянка» этапов на линейке]
Требуется реализовать мигание заданное число раз для каждого этапа, используя короткие функции и отложенные вызовы.

\textbf{Ожидаемый результат:} корректный световой отчёт по морзянке для каждого этапа.
\end{taskbox}

\begin{taskbox}[Траектория «лепесток»]
Требуется реализовать формулу \(r = r_0 \cdot (1 + 0.2 \sin(k \cdot \theta))\):
\begin{itemize}
    \item при \(\theta = \text{angle} \cdot \pi / 180\);
    \item выбрать \(k \ge 3\) и обновлять точку параметрически.
\end{itemize}

\textbf{Ожидаемый результат:} визуально различимый «лепесток» по параметрической траектории.
\end{taskbox}

\begin{taskbox}[Проверка точности навигации]
Требуется обход 8 точек окружности, измерение интервалов \texttt{dt} между \texttt{Ev.POINT\_REACHED} с использованием \texttt{os.clock()}.

\textbf{Ожидаемый результат:} получены интервалы времени, оценено постоянство периода между достижениями точек.
\end{taskbox}

\begin{taskbox}[Миссия «периметр» вокруг базы]
Требуется определить набор точек по квадрату вокруг \texttt{(0,0)}, включить возврат в исходную точку и выполнить посадку.

\textbf{Ожидаемый результат:} полный обход периметра и посадка в исходной точке.
\end{taskbox}

\begin{taskbox}[«Окно безопасности» для \texttt{yaw}]
Требуется ограничить изменение \texttt{yaw} на тик не более \texttt{maxYawRate * dt} и обновлять точку только при соблюдении ограничения.

\textbf{Ожидаемый результат:} курс изменяется не быстрее установленной скорости, миссия сохраняет устойчивость.
\end{taskbox}

\begin{taskbox}[Демонстрационный сценарий (RC старт, \(\approx\) 2 минуты, посадка)]
Требуется запуск по RC, полёт по окружности около двух минут, посадка, световой отчёт этапов, отсутствие блокировок и проверка безопасного завершения.

\textbf{Ожидаемый результат:} демонстрационный полёт по окружности, корректная посадка и полная индикация этапов.
\end{taskbox}
